\clearpage
\SourceLeaf{604}
\NoSourceHead
\begin{center}\Large\bfseries BERICHTIGUNGEN.\end{center}
\begin{center}\rule{.1\linewidth}{.4pt}\end{center}
\medskip
S.\ 7. Die Überschrift soll lauten: Vorrede. Mathematische Abhandlungen besonders aus dem Gebiete der Höheren Arithmetik und der Elliptischen Functionen von Dr.\ G.\ \textsc{Eisenstein}. Berlin 1847.

S.\ 110, Zeile 1, 2. Die Urschrift des Briefes von \textsc{Gauss} an \textsc{Drobisch} vom 14.\ August 1834 befindet sich im \textsc{Gauss}archiv.

S.\ 148, in der Formel für $s_4$ (letzte Gleichung des art.\ [6.]) muß der Nenner des zweiten Bruches lauten $1+20s^4-26s^8+20s^{12}+s^{16}$.

S.\ 150, Gleichung (3) muß lauten $AB=\dfrac{\pi}{4}$.

S.\ 171, Zeile 9 ist das Komma hinter \textls[120]{Scheda} zu streichen.

»\quad »\,, Zeile 20, 21 in der Klammer ist zu lesen: es folgen nämlich auf die dreizehnte Decimalstelle 0 noch die Stellen 73164.

S.\ 188, wo die Nummer [24] zweimal auftritt, ist das erste Mal [23] zu lesen.

S.\ 192, letzte Zeile ist statt $\left(\dfrac{1}{2}\right)^2$ zu lesen $\left(\dfrac{1}{2}\right)^3$.

S.\ 251, zweite Fußnote ist statt S.\ 247 zu lesen S.\ 249.

S.\ 254, Gleichung (19) ist statt $\left(\dfrac{1\mathbin{.}3}{2\mathbin{.}4}\right)^2k^2$ zu lesen $\left(\dfrac{1\mathbin{.}3}{2\mathbin{.}4}\right)^2k^4$.

S.\ 379, dritte Fußnote ist statt \emph{novo} zu lesen \emph{nova}.

S.\ 389, in der Fußnote ist statt $\dfrac{1}{2}\pi\left(k+\dfrac{1}{2}\right)$ zu lesen $\dfrac{1}{2}\pi\left(k+\dfrac{1}{4}\right)$.

S.\ 444, vorletzte Textzeile und

S.\ 445, erste Fußnote ist statt 27.\ Januar zu lesen 24.\ Januar.

S.\ 509, fünfte Zeile der Bemerkung zu der Nr.\ [49] ist statt *) zu lesen **).

S.\ 535, Zeile 11 ist hinter »$=0$« ein Komma zu setzen.
\ClosingRule
